\scanpage{57}
\uline{\unsure{Laius|Laïus}}\hfill 6.12.1966
\diagram{\begin{tikzcd}[ampersand replacement=\&,row sep=small]
X \arrow[d,"f"']\\ S
\end{tikzcd}}

Soit $F$ un Module quasi-cohérent stratifié sur $X/S$. On a vu qu'il définit un module de $X_{\mathrm{strat}}$, et par là un faisceau cosimplicial \inserted{$F^{\unsure{+|*}}=C(F)$} sur $X_{\mathrm{zar}}$, par
\[
C^n(F)=F^{(n)}=\varprojlim_i F_{\Delta^n_{X/S}(i)}
\]
[peu importe laquelle des $n+1$ projections $\Delta^{\unsure{-|n}}_{X/S}(i)\longrightarrow X$ on prend pour définir l'image inverse, il y a des isom canoniques entre eux, précisément à cause de la stratification de $F$].

\unsure{On|Cela} a
\[
\boxed{H^*(X_{\mathrm{str}},F)\simeq\mathbb H^*(X_{\mathrm{zar}},\mathbb C^*(F))}
\]

On voit donc deux théories cohomologiques ; \deleted{\unsure{qui|par}} \inserted{\deleted{dont}} \deleted{aucune ne semble \unsure{pas|devoir} à première vue} \deleted{semblent assez \unsure{distinctes} \unsure{qu'il|gén\textsuperscript{t}} \unsure{ne} pas contenir l'autre} :

a) Pour un $F$ quasi-cohérent sur $X$, on forme les $H^i(X_{\mathrm{zar}},F)$.

b) Si $F$ est \uline{stratifié}, on forme les $H^i(X_{\mathrm{strat}},F)$, qui s'explicitent par la formule plus haut.

À première vue, le lien entre ces deux théories cohomologiques semble donné exclusivement par
% archive sheet number: 56

